\begin{frame}{세 번째 교훈: 라벨의 형태가 모델의 형태를 결정한다}
  \begin{itemize}
    \item 2012년의 패러다임 전환은 ``\textbf{사진 + 범주 라벨}''이라는
      \kb{지도학습}의 가장 전형적인 형태 위에서 일어났다.
    \item 정답이 없는 데이터만 있을 때(비지도학습, Week 14$\sim$15),
      보상 신호만 있을 때(강화학습, ML2)에는 같은 GPU·데이터 축도
      중요하지만,
  \end{itemize}
  \vskip0.4em
  \kq{``어떤 신호로 학습하는가''에 따라 최적의 모델 가족 자체가 달라진다.}
  \vskip0.6em
  \begin{itemize}
    \item 1.2절에서 ``\textbf{데이터의 형태가 문제를 결정한다}''고 할 때
      뜻하는 바가 정확히 이것.
    \item 부수 효과: ``우리 데이터가 작으니 신경망은 아직 이르다'' 같은
      실무 판단이 왜 흔한지도 자연스럽게 이해된다.
  \end{itemize}
\end{frame}
